\documentclass[11pt,a4paper]{article}

\usepackage[margin=1in]{geometry}
\usepackage{amsmath}
\usepackage{amssymb}
\usepackage{booktabs}
\usepackage{hyperref}

\hypersetup{colorlinks=true, linkcolor=black, urlcolor=blue, citecolor=black}

\newcommand{\cacose}{\texttt{CaCoSE}}
\newcommand{\caef}{\textit{CaEF}}

\title{\vspace{-2em}CaCoSE Phase 1: Implementation Specification\\
\large A from-the-paper reimplementation of cross-attentive cohesive subgraph embedding}
\author{GNN\_CaCoSE project}
\date{26 August 2026}

\begin{document}
\maketitle
\vspace{-1.5em}

\begin{abstract}
\noindent This document fixes the Phase~1 design of a reimplementation of \cacose{} (Hossain,
Islam, Chebbah, Fanning and Akbas, \textit{Cross-attentive Cohesive Subgraph Embedding to Mitigate
Oversquashing in GNNs}, arXiv:2603.27529v3) before any code is written. No reference implementation
of the method is publicly available, so every construction here is derived from the paper text,
its Figure~2 worked example, and Algorithm~1 in its appendix. Where the paper underdetermines an
implementation choice, this document records what the paper says, what we implement, and whether
the reading is settled or still open. Phase~1 succeeds when the package reproduces the paper's reported
accuracy on Cora, Chameleon and MUTAG to within a stated tolerance.
\end{abstract}

\section{Scope}

Phase~1 delivers a tested Python package, \texttt{cacose}, built on PyTorch Geometric, and a
reproduction of three of the paper's twelve reported results. The three datasets span the regimes
the method claims to serve: Cora (homophilic node classification), Chameleon (heterophilic node
classification, where the paper claims its largest margin) and MUTAG (graph classification).

\begin{table}[!htbp]
\centering
\begin{tabular}{llrr}
\toprule
Dataset & Task & Paper accuracy & Phase 1 accepts at \\
\midrule
Cora      & Node classification  & 85.00 & $\geq$ 83.5 \\
Chameleon & Node classification  & 68.99 & $\geq$ 66.5 \\
MUTAG     & Graph classification & 76.99 & $\geq$ 74.0 \\
\bottomrule
\end{tabular}
\caption{Reproduction targets, mean over 10 seeds (paper Tables 1 and 2).}
\end{table}

Out of scope for Phase~1: application of the method to a new domain (Phase~2) and an interactive
attention visualiser (Phase~3). Both later phases consume this package rather than modify it, which
constrains two interfaces described in Section~\ref{sec:hooks}.

\section{Method}

A graph is $G = (V, E, X)$ with $|V| = N_v$ nodes and features $X \in \mathbb{R}^{N_v \times d}$.
The method is a deterministic graph transform followed by a trainable network. Separating the two
matters for engineering: the transform runs once per dataset on CPU and is cached, and the network
never modifies it.

\subsection{Module 1: cohesive subgraph decomposition}

Let $\mathrm{core}(v)$ be the node core number, and $G_k$ the node-induced subgraph on
$\{v : \mathrm{core}(v) \geq k\}$. The paper's edge score (its Definition~2) is
\begin{equation}
C(u,v) = \max\{k \mid (u,v) \in G_k\}.
\end{equation}
Because $(u,v) \in G_k$ exactly when both endpoints survive to level $k$, this reduces to
\begin{equation}
C(u,v) = \min\bigl(\mathrm{core}(u), \mathrm{core}(v)\bigr),
\label{eq:score}
\end{equation}
which we implement directly from a single call to \texttt{networkx.core\_number}. Materialising each
$G_k$ to compute scores, as a literal reading of the definition suggests, is unnecessary.

Closure-aware edge filtration (\caef{}, the paper's Theorem~1) then demotes edges that carry no
triadic support. For an edge with score $k \geq \delta$,
\begin{equation}
S(u,v) = \bigl|N_k(u) \cap N_k(v)\bigr|, \qquad
S(u,v) = 0 \;\Longrightarrow\; C(u,v) \leftarrow k - 1,
\label{eq:caef}
\end{equation}
where $N_k(\cdot)$ is the neighbourhood \emph{inside $G_k$}. The restriction to $G_k$ rather than
$G$ is the single most consequential reading in this document, and Section~\ref{sec:example}
shows that the paper's own Figure~2 is only consistent under it. Appendix~A.1 of the paper supplies
the justification for the rule itself: an edge with $N(u) \cap N(v) = \emptyset$ has non-positive
Ollivier-Ricci curvature, so \caef{} removes exactly the edges that behave as bottlenecks.

Subgraphs are then the score classes,
\begin{equation}
S = \{S_k\}_{k=1}^{k_{\max}}, \qquad S_k = (V_k, E_k), \qquad E_k = \{(u,v) \in E \mid C(u,v) = k\}.
\label{eq:partition}
\end{equation}
This partitions the \emph{edge} set. A node may belong to many subgraphs, an edge to exactly one,
and $S_k \neq G_k$ in general. The number of subgraphs is $k_{\max}$, the graph degeneracy.

\subsection{Module 2: graph learning}

Each subgraph is embedded independently, with its own parameters:
\begin{equation}
H_k = \mathrm{GCN}_k(S_k), \qquad Z_k = \mathrm{SAGPool}_k(S_k, H_k).
\end{equation}
$H_k \in \mathbb{R}^{|V_k| \times h}$ holds node embeddings local to the subgraph; $Z_k \in
\mathbb{R}^{d_S}$ is a single pooled vector summarising it. Self-attention graph pooling selects the
top $\lceil \mathrm{PR} \cdot |V_k| \rceil$ nodes by a GCN-computed attention score before readout,
which is where the paper's claim of filtering task-irrelevant neighbours is realised.

\subsection{Module 3: cross-subgraph attention and feature combination}

Pooling discards the relationships \emph{between} regions, so the subgraph embeddings are stacked
into $Z_S \in \mathbb{R}^{N_S \times d_S}$ and passed through scaled dot-product self-attention:
\begin{equation}
\mathrm{Attn}_S = \mathrm{softmax}\!\left(\frac{QK^{\top}}{\sqrt{d_S}}\right), \qquad
Z_S^{\mathrm{attn}} = \mathrm{Attn}_S V,
\label{eq:attn}
\end{equation}
with $Q = Z_S W_Q$, $K = Z_S W_K$, $V = Z_S W_V$. Since $N_S = k_{\max}$ is small, this costs little.
The graph representation is the mean over subgraphs,
\begin{equation}
Z_G = \mu\bigl([Z_1^{\mathrm{attn}}, \ldots, Z_{k_{\max}}^{\mathrm{attn}}]\bigr).
\end{equation}
Node representations concatenate each node's local embedding with its subgraph's attended embedding
and sum over every subgraph the node appears in:
\begin{equation}
h_v^{(k)} = h_v \,\|\, Z_k^{\mathrm{attn}} \;\;(v \in S_k), \qquad
z_v = \sum_{k \,:\, v \in S_k} h_v^{(k)}.
\label{eq:merge}
\end{equation}
The final node representation therefore has width $h + d_S$. Nodes isolated by the decomposition
receive the zero vector; we count and log them rather than letting them pass silently.

\subsection{Module 4: prediction}

A two-layer perceptron maps $z_v$ to node labels or $Z_G$ to graph labels, trained with cross
entropy.

\subsection{Procedure}

\noindent\hrulefill
\begin{enumerate}\setlength{\itemsep}{0pt}
\item Strip self-loops, symmetrise $E$. Compute $\mathrm{core}(v)$ for all $v$.
\item Score every edge by Equation~\eqref{eq:score}.
\item For each distinct $k \geq \delta$: build $G_k$ once, and apply Equation~\eqref{eq:caef} to the
      edges scoring $k$.
\item Form $\{S_k\}$ by Equation~\eqref{eq:partition}; emit each as a local edge index plus a
      local-to-global node map. Cache the result.
\item For each $S_k$: $H_k \leftarrow \mathrm{GCN}_k(S_k)$, $Z_k \leftarrow
      \mathrm{SAGPool}_k(S_k, H_k)$.
\item $Z_S^{\mathrm{attn}} \leftarrow \mathrm{Attention}(Z_S)$ by Equation~\eqref{eq:attn}; retain
      the attention matrix.
\item $Z_G \leftarrow \mu(Z_S^{\mathrm{attn}})$; accumulate $z_v$ by Equation~\eqref{eq:merge}.
\item Return $(Z_v, Z_G)$ and the attention matrix.
\end{enumerate}
\noindent\hrulefill

\medskip
Steps 1--4 are deterministic and cached per $(\textrm{dataset}, \delta, \textrm{caef\_mode})$.
Step 5 is parallel across $k$.

\subsection{Interfaces}
\label{sec:interfaces}

Two structures are fixed by this specification rather than left to the implementation, because
later phases of the project bind to them directly.

\begin{small}
\begin{verbatim}
Subgraph(k: int, edge_index: Tensor, node_map: Tensor)
Decomposition(subgraphs: list[Subgraph], num_nodes: int, kmax: int)
    cache key: (dataset, decomposer_id, param_hash)

forward(...) -> (logits, aux)
    aux = {Z_S, Z_S_attn, attn_weights, subgraph_ids}
\end{verbatim}
\end{small}

\noindent
\texttt{node\_map} carries local-to-global node indices, so a subgraph's tensors can be interpreted
without reference to the graph it came from; \texttt{edge\_index} is in local indexing. The cached
\texttt{Decomposition} is the entry point for applying the method to a new domain. In \texttt{aux},
\texttt{attn\_weights} is retained per head rather than averaged, which is what an attention
visualiser needs. Both structures are covered by tests asserting their fields and shapes.

The cache key generalises what an earlier draft fixed as $(\textrm{dataset}, \delta,
\textrm{caef\_mode})$. That form names parameters belonging to one decomposer; since the paper's
own Figure~5 compares against Louvain and Metis partitions, the key must identify the decomposer
too. For \texttt{KCoreCaEF} the parameter hash covers exactly $\delta$ and the \caef{} mode, so the
two forms agree.

\section{Assumptions and ambiguity resolutions}
\label{sec:ambiguity}

The table below lists every point where the paper underdetermines the implementation. A status of
Resolved means the paper's own text or figures settle the question; Open means we proceed on a
working assumption. Open items are candidates for correspondence with the first author.

\begin{center}\small
\begin{tabular}{@{}p{0.5cm}p{3.1cm}p{4.3cm}p{4.6cm}p{1.6cm}@{}}
\toprule
\# & Question & What we implement & Reason & Status \\
\midrule
1 & \caef{} decrement: once, or cascade to $k-2$? & Single decrement, configurable & Algorithm~1 lines 2--3 show one assignment, no loop & Open \\
2 & Triadic support in $G_k$ or in $G$? & Inside $G_k$ & Only reading consistent with Figure~2; see Section~\ref{sec:example} & \textbf{Resolved} \\
3 & $\mathrm{GCN}_k$ weights separate or shared? & Separate, flag to share & Subscript $k$ in Equation~(6) & Open \\
4 & SAGPool readout function & mean $\|$ max, configurable & Standard SAGPool-g practice; paper says only ``READOUT'' & Open \\
5 & Residual/LayerNorm on attention? & Absent, flag to enable & Equation~(8) is stated plainly, without either & Open \\
6a & Cora split fixed or per seed? & Re-randomised per seed & Paper reports a mean over 10 seeds & Open \\
6b & Cora counts 1208/500/500 leave 500 nodes unused & Literal 1208/500/500 & 1208+500+1000 = 2708 exactly, so ``500'' may be a typo for the standard split; we follow the text & Open \\
7 & Chameleon: which version and split? & Geom-GCN preprocessed graph, our own stratified 48/32/20 per seed & Paper states 48/32/20; the packaged Geom-GCN masks are 60/20/20 & \textbf{Resolved} \\
8 & Meaning of ``1-hot encoding'' features & Native dataset features & MUTAG already carries one-hot atom types; the remark most plausibly targets featureless GC datasets & Open \\
9 & GCN depth per subgraph & Two layers & Unstated; two is the default in the cited GCN work & Open \\
10 & Is $d_S$ equal to the GCN hidden width? & Both 128 & Paper gives one hidden dimension, 128 & Open \\
11 & Dropout rate --- never stated & 0.5 for NC, 0 for GC & Measured: at 0.5 the MUTAG eval loss diverges from training loss ($0.47 \rightarrow 1.03$) and predictions collapse to one class, because dropout perturbs the features SAGPool selects nodes from & Open \\
\bottomrule
\end{tabular}

\vspace{0.4em}
{\small Ambiguity log. Items 2 and 7 are settled; the rest are working assumptions.
 Item 11 is not in the paper at all --- it is a hyperparameter the text never mentions.}
\end{center}

\section{Worked example}
\label{sec:example}

The paper's Figure~2 provides a ten-node example that pins down the decomposition exactly. It
consists of a $K_5$ on $v_1 \ldots v_5$, a triangle $v_4 v_6 v_7$, and a $K_4$ on
$v_7 \ldots v_{10}$. Core numbers are $\mathrm{core}(v_1 \ldots v_5) = 4$,
$\mathrm{core}(v_7 \ldots v_{10}) = 3$ and $\mathrm{core}(v_6) = 2$.

\begin{center}\small
\begin{tabular}{@{}llccl@{}}
\toprule
Edges & Count & Raw $C(u,v)$ & $S(u,v)$ at that $k$ & Final $C(u,v)$ \\
\midrule
$K_5$ interior                & 10 & 4 & 3 & 4 \\
$K_4$ interior                &  6 & 3 & 2 & 3 \\
$(v_4,v_7)$                   &  1 & 3 & \textbf{0} & \textbf{2} \\
$(v_4,v_6)$, $(v_6,v_7)$      &  2 & 2 & not checked ($k < \delta$) & 2 \\
\bottomrule
\end{tabular}

\vspace{0.4em}
{\small Edge scores on the Figure~2 graph with $\delta = 3$.}
\end{center}

The decisive row is $(v_4, v_7)$. In the full graph the two endpoints share the neighbour $v_6$, so
support computed on $G$ would be~1 and the edge would survive at level~3. But $\mathrm{core}(v_6) =
2$, so $v_6$ is absent from $G_3$, support inside $G_3$ is~0, and the edge is demoted to level~2.
The paper states that this edge ``has no support'' at $k=3$, which is true only under the
$G_k$-restricted reading. This settles ambiguity~2.

The resulting partition is $\{S_2, S_3, S_4\}$ with $|E(S_2)| = 3$, $|E(S_3)| = 6$ and
$|E(S_4)| = 10$, and $v_4$ belongs to $S_2$ and $S_4$ but not $S_3$ --- matching the paper's
statement that $z_{v_4} = h_{v_4(2)} + h_{v_4(4)}$.

This example is the primary unit-test fixture. Encoding it once, here and in the test suite, keeps
the specification and the implementation from drifting apart.

\section{Experimental protocol}

Hyperparameters follow the paper's Section~4.1 without modification: Adam at learning rate
$2.5\times10^{-3}$, weight decay $1\times10^{-4}$, hidden width 128, pooling ratio 0.5, $\delta = 3$,
two attention heads for node classification and one for graph classification. Node classification
trains for at most 250 epochs with early stopping after 50 epochs without validation improvement;
graph classification uses 100 and 25. Every configuration is run over 10 seeds and reported as a
mean.

Splits follow the paper as written: Cora 1208/500/500, Chameleon 48/32/20 stratified, MUTAG
80/10/10 stratified, all re-drawn per seed. Learning rate and hidden width are treated as fixed by
the paper and are never tuned; if a target is missed, only the items left open in the
ambiguity table of Section~\ref{sec:ambiguity} are varied, in the order 1, 4, 5, 3, then dropout.

\section{Validation and risks}
\label{sec:hooks}

Correctness rests on the decomposition, which is deterministic and therefore fully testable. Beyond
the Figure~2 fixture, we assert three properties over random graphs: every input edge lands in
exactly one $S_k$; the union of the node maps is the set of non-isolated nodes; and for every edge
whose raw score $k$ is at least $\delta$, zero support at level $k$ holds if and only if its final
score is $k-1$. The last statement is the correct form of the invariant --- the simpler claim that
every retained edge at level $\geq \delta$ has positive support is false, because a demoted edge can
land at a level that is still $\geq \delta$ and is never rechecked under a single decrement.

The two interfaces of Section~\ref{sec:interfaces} are fixed now because later phases bind to them:
the cached \texttt{Decomposition} is the Phase~2 entry point, and the \texttt{aux} dictionary is what
the Phase~3 visualiser reads. Changing either after Phase~1 means changing both consumers.

One risk deserves stating in advance. Chameleon has a degeneracy near 63, judging from the $k \in
\{61, 62, 63\}$ values used in the paper's own pilot study, so the model instantiates roughly 63
separate GCN and pooling branches. With a 2325-dimensional input and unshared weights that is on the
order of 19 million parameters in the first layers alone, fitted to 2277 nodes. If Chameleon falls
short of its target, shared weights (ambiguity~3) is the first variant to try. Cora and MUTAG, with
degeneracies of 4 and roughly 3, are unaffected.

\end{document}
